\section*{Twice 2 Bits in 3 Bits Instructions}

\subsection*{Objective}
Design a combinational read circuit that decodes a 3-bit write-once memory snapshot into the currently represented 2-bit value.

\subsection*{Problem Background}
This puzzle models a memory mechanism with \textbf{write-once bits}. A bit may change from 0 to 1, but it can never return to 0.

The broader challenge is to represent one 2-bit value, and later represent another 2-bit value, while reusing as few physical bits as possible. During the process, a reader may inspect memory at any moment and must still recover the value currently intended by the writer.

In this specific puzzle, the encoding side is fixed to a 3-bit memory state (x1, x2, x3). Your task is to build only the \textbf{reader}: given those three bits, output the correct decoded 2-bit value.

\subsection*{What You Must Build}
Create a logic circuit that maps each possible 3-bit memory state to the correct 2-bit decoded value.

The circuit must be purely combinational with respect to the three inputs:
\begin{itemize}
  \item Inputs: x1, x2, x3
  \item Outputs: out1, out0
\end{itemize}

\subsection*{Output Format}
The decoded value is represented in binary using two output wires:
\begin{itemize}
  \item out1: most-significant bit (MSB)
  \item out0: least-significant bit (LSB)
\end{itemize}

So a decoded value v in {0,1,2,3} must be emitted as:
\begin{itemize}
  \item 0 -> 00
  \item 1 -> 01
  \item 2 -> 10
  \item 3 -> 11
\end{itemize}

\subsection*{Design Intent}
This is a reasoning puzzle about information reuse under one-way memory updates. You are expected to infer the read logic from the write-once model and the requirement that the represented value must remain decodable at read time.

To make the target mapping unique, use the following two anchor constraints:
\begin{itemize}
  \item \textbf{Canonical one-hot anchors:} 000 decodes to value 0, 100 decodes to 1, 010 decodes to 2, and 001 decodes to 3.
  \item \textbf{Complement invariance:} complementing all three memory bits does not change the decoded value.
\end{itemize}

\subsection*{Requirements}
\begin{itemize}
  \item Implement a decoder for inputs x1, x2, x3.
  \item Produce outputs out1, out0 as a valid 2-bit value encoding.
  \item The circuit must handle all 8 input combinations correctly.
\end{itemize}

\subsection*{Evaluation}
A submission is accepted only if it matches the hidden reference behavior for every tested input state.
